% GENERATED FILE. Edit canonical lessons, then run npm run generate:books.
% canonical-source-hash: fnv1a64:5308a4cf560574d1
% canonical-lessons: MR-C519-vijeri, MR-C519-jhumbar, MR-C519-ujed, MR-C519-andhar, MR-C519-ujalne, MR-R519-review-512-519

\chapter{Light at Home}
\label{ch:mr-light-at-home}
\hlchaptermodality{519}

\begin{chapteropening}
\textbf{By the end of this chapter:} \emph{I can talk about light and darkness at home.}

Complete the last lesson of chapter 519: I can talk about light and darkness at home.
\end{chapteropening}

\section[\texorpdfstring{\mr{विजेरी} (vijerī)}{विजेरी (vijerī)}]{\mr{विजेरी} (vijerī) — a torch; a flashlight}

\label{lesson:MR-C519-vijeri}



\begin{quote}
Before the new one: say the Marathi for automatic, then the Marathi for to spray.
\end{quote}

\subsection*{You'll want to know: \mr{विजेरी}}
\textbf{\mr{विजेरी}} — \emph{vijerī} — \textquotedblleft{}a torch; a flashlight\textquotedblright{}.

\begin{quote}
\mr{वीज} \mr{गेली}, \mr{म्हणून} \mr{मी} \mr{विजेरी} \mr{लावली}.
\end{quote}

\begin{quote}
\emph{vīj gelī, mhaṇūn mī vijerī lāvlī.}
\end{quote}

\begin{quote}
The power went, so I switched on a torch.
\end{quote}

Feminine.

\begin{cousinweb}
\textbf{\mr{विजेरी}} (\emph{vijerī}) is built on \mr{वीज}, \textquotedblleft{}lightning, electricity\textquotedblright{}, inherited from Sanskrit \emph{vidyut} (CDIAL 11742), with a Marathi ending.
\end{cousinweb}

\begin{grammarlens}[title={What you've built}]
One more word for describing what you lived and felt.
\end{grammarlens}

\subsection*{Your turn}
\begin{itemize}
\raggedright
  \item \emph{Say it:} \emph{vijerī}
  \item \emph{Say it:} \emph{vijerī}, once more
  \item \emph{Say it:} \emph{vīj gelī, mhaṇūn mī vijerī lāvlī.} — The power went, so I switched on a torch.
  \item \emph{Recall:} say the Marathi for automatic, then the Marathi for to spray, then say \emph{vijerī} again
\end{itemize}

\begin{tcolorbox}[breakable,colback=teal!4,colframe=teal!35!black,title={Before you move on}]
What does \mr{विजेरी} mean? (\textquotedblleft{}A torch; a flashlight\textquotedblright{}.) Say it once more.
\end{tcolorbox}
\section[\texorpdfstring{\mr{झुंबर} (jhumbar)}{झुंबर (jhumbar)}]{\mr{झुंबर} (jhumbar) — a chandelier}

\label{lesson:MR-C519-jhumbar}



\begin{quote}
Before the new one: say the Marathi for to spray, then the Marathi for a torch.
\end{quote}

\subsection*{You'll want to know: \mr{झुंबर}}
\textbf{\mr{झुंबर}} — \emph{jhumbar} — \textquotedblleft{}a chandelier\textquotedblright{}.

\begin{quote}
\mr{जुन्या} \mr{घरात} \mr{काचेचं} \mr{मोठं} \mr{झुंबर} \mr{होतं}.
\end{quote}

\begin{quote}
\emph{junyā gharāt kācecaṁ moṭhaṁ jhumbar hotaṁ.}
\end{quote}

\begin{quote}
There was a big glass chandelier in the old house.
\end{quote}

Neuter.

\begin{cousinweb}
\textbf{\mr{झुंबर}} (\emph{jhumbar}) names a hanging cluster of lamps; its origin is not settled.
\end{cousinweb}

\begin{grammarlens}[title={What you've built}]
One more word for describing what you lived and felt.
\end{grammarlens}

\subsection*{Your turn}
\begin{itemize}
\raggedright
  \item \emph{Say it:} \emph{jhumbar}
  \item \emph{Say it:} \emph{jhumbar}, once more
  \item \emph{Say it:} \emph{junyā gharāt kācecaṁ moṭhaṁ jhumbar hotaṁ.} — There was a big glass chandelier in the old house.
  \item \emph{Recall:} say the Marathi for to spray, then the Marathi for a torch, then say \emph{jhumbar} again
\end{itemize}

\begin{tcolorbox}[breakable,colback=teal!4,colframe=teal!35!black,title={Before you move on}]
What does \mr{झुंबर} mean? (\textquotedblleft{}A chandelier\textquotedblright{}.) Say it once more.
\end{tcolorbox}
\section[\texorpdfstring{\mr{उजेड} (ujeḍ)}{उजेड (ujeḍ)}]{\mr{उजेड} (ujeḍ) — daylight; light}

\label{lesson:MR-C519-ujed}



\begin{quote}
Before the new one: say the Marathi for a torch, then the Marathi for a chandelier.
\end{quote}

\subsection*{You'll want to know: \mr{उजेड}}
\textbf{\mr{उजेड}} — \emph{ujeḍ} — \textquotedblleft{}daylight; light\textquotedblright{}.

\begin{quote}
\mr{सकाळी} \mr{खोलीत} \mr{छान} \mr{उजेड} \mr{येतो}.
\end{quote}

\begin{quote}
\emph{sakāḷī kholīt chān ujeḍ yeto.}
\end{quote}

\begin{quote}
In the morning lovely light comes into the room.
\end{quote}

Masculine.

\mr{उजेड} \mr{येणे}, light comes in; \mr{उजेड} \mr{पडणे}, light falls.

\begin{cousinweb}
\textbf{\mr{उजेड}} (\emph{ujeḍ}) is a tadbhava: CDIAL lists Marathi \emph{ujeḍ}, \textquotedblleft{}light\textquotedblright{}, under Sanskrit \emph{uddyōta}, \textquotedblleft{}flashing, light\textquotedblright{} (CDIAL 2000).
\end{cousinweb}

\begin{grammarlens}[title={What you've built}]
One more word for describing what you lived and felt.
\end{grammarlens}

\subsection*{Your turn}
\begin{itemize}
\raggedright
  \item \emph{Say it:} \emph{ujeḍ}
  \item \emph{Say it:} \emph{ujeḍ}, once more
  \item \emph{Say it:} \emph{sakāḷī kholīt chān ujeḍ yeto.} — In the morning lovely light comes into the room.
  \item \emph{Recall:} say the Marathi for a torch, then the Marathi for a chandelier, then say \emph{ujeḍ} again
\end{itemize}

\begin{tcolorbox}[breakable,colback=teal!4,colframe=teal!35!black,title={Before you move on}]
What does \mr{उजेड} mean? (\textquotedblleft{}Daylight; light\textquotedblright{}.) Say it once more.
\end{tcolorbox}
\section[\texorpdfstring{\mr{अंधार} (andhār)}{अंधार (andhār)}]{\mr{अंधार} (andhār) — darkness}

\label{lesson:MR-C519-andhar}



\begin{quote}
Before the new one: say the Marathi for a chandelier, then the Marathi for daylight.
\end{quote}

\subsection*{You'll want to know: \mr{अंधार}}
\textbf{\mr{अंधार}} — \emph{andhār} — \textquotedblleft{}darkness\textquotedblright{}.

\begin{quote}
\mr{रात्री} \mr{बागेत} \mr{खूप} \mr{अंधार} \mr{असतो}.
\end{quote}

\begin{quote}
\emph{rātrī bāget khūp andhār asto.}
\end{quote}

\begin{quote}
At night it is very dark in the garden.
\end{quote}

Masculine.

\begin{cousinweb}
\textbf{\mr{अंधार}} (\emph{andhār}) is a tadbhava, inherited from Sanskrit \emph{andhakāra}, \textquotedblleft{}darkness\textquotedblright{} (CDIAL 386).
\end{cousinweb}

\begin{grammarlens}[title={What you've built}]
One more word for describing what you lived and felt.
\end{grammarlens}

\subsection*{Your turn}
\begin{itemize}
\raggedright
  \item \emph{Say it:} \emph{andhār}
  \item \emph{Say it:} \emph{andhār}, once more
  \item \emph{Say it:} \emph{rātrī bāget khūp andhār asto.} — At night it is very dark in the garden.
  \item \emph{Recall:} say the Marathi for a chandelier, then the Marathi for daylight, then say \emph{andhār} again
\end{itemize}

\begin{tcolorbox}[breakable,colback=teal!4,colframe=teal!35!black,title={Before you move on}]
What does \mr{अंधार} mean? (\textquotedblleft{}Darkness\textquotedblright{}.) Say it once more.
\end{tcolorbox}
\section[\texorpdfstring{\mr{उजळणे} (ujaḷṇe)}{उजळणे (ujaḷṇe)}]{\mr{उजळणे} (ujaḷṇe) — to light up; to brighten}

\label{lesson:MR-C519-ujalne}



\begin{quote}
Before the new one: say the Marathi for daylight, then the Marathi for darkness.
\end{quote}

\subsection*{You'll want to know: \mr{उजळणे}}
\textbf{\mr{उजळणे}} — \emph{ujaḷṇe} — \textquotedblleft{}to light up; to brighten\textquotedblright{}.

\begin{quote}
\mr{दिवे} \mr{लागले} \mr{आणि} \mr{सगळा} \mr{रस्ता} \mr{उजळला}.
\end{quote}

\begin{quote}
\emph{dive lāgle āṇi sagḷā rastā ujaḷlā.}
\end{quote}

\begin{quote}
The lamps came on and the whole street lit up.
\end{quote}

Past (he …): \mr{तो} \mr{रस्ता} \mr{उजळला}, \emph{to rastā ujaḷlā}.

\begin{cousinweb}
\textbf{\mr{उजळणे}} (\emph{ujaḷṇe}) is a tadbhava: CDIAL lists Marathi \emph{ujaḷṇẽ} under Sanskrit \emph{ujjvalati}, \textquotedblleft{}blazes up, shines\textquotedblright{} (CDIAL 1671).
\end{cousinweb}

\begin{grammarlens}[title={What you've built}]
One more word for describing what you lived and felt.
\end{grammarlens}

\subsection*{Your turn}
\begin{itemize}
\raggedright
  \item \emph{Say it:} \emph{ujaḷṇe}
  \item \emph{Say it:} \emph{ujaḷṇe}, once more
  \item \emph{Say it:} \emph{dive lāgle āṇi sagḷā rastā ujaḷlā.} — The lamps came on and the whole street lit up.
  \item \emph{Recall:} say the Marathi for daylight, then the Marathi for darkness, then say \emph{ujaḷṇe} again
\end{itemize}

\begin{tcolorbox}[breakable,colback=teal!4,colframe=teal!35!black,title={Before you move on}]
What does \mr{उजळणे} mean? (\textquotedblleft{}To light up; to brighten\textquotedblright{}.) Say it once more.
\end{tcolorbox}
\section[(dialogue)]{Looking back: from Made in a Factory to Light at Home}

\label{lesson:MR-R519-review-512-519}



\begin{quote}
Say the Marathi for a factory and production.
\end{quote}

\subsection*{Your turn}
Say these aloud:

\begin{itemize}
\raggedright
  \item a factory, production, goods, a supply, to sell well
  \item a workshop, a command, a department, to present, to host a programme
  \item a task, the day's business, a shift, to get done, to struggle
  \item an engineer, a consultant, a translator, an accountant, to wrap up
  \item a computer, a mobile phone, a telephone, a video clip, to use
  \item a password, the internet, a message, a link, to store
  \item a machine, a device, a robot, automatic, to spray
  \item a torch, a chandelier, daylight, darkness, to light up
\end{itemize}

\begin{tcolorbox}[breakable,colback=teal!4,colframe=teal!35!black,title={Before you move on}]
A factory? (\textbf{\mr{कारखाना}}, \emph{kārkhānā}.) Production? (\textbf{\mr{उत्पादन}}, \emph{utpādan}.) Goods? (\textbf{\mr{माल}}, \emph{māl}.) A supply? (\textbf{\mr{पुरवठा}}, \emph{puravṭhā}.) To sell well? (\textbf{\mr{खपणे}}, \emph{khapṇe}.) A workshop? (\textbf{\mr{कार्यशाळा}}, \emph{kāryaśāḷā}.) A command? (\textbf{\mr{आदेश}}, \emph{ādeś}.) A department? (\textbf{\mr{विभाग}}, \emph{vibhāg}.) To present? (\textbf{\mr{सादर} \mr{करणे}}, \emph{sādar karṇe}.) To host a programme? (\textbf{\mr{सूत्रसंचालन} \mr{करणे}}, \emph{sūtrasañcālan karṇe}.) A task? (\textbf{\mr{कार्य}}, \emph{kārya}.) The day's business? (\textbf{\mr{कामकाज}}, \emph{kāmkāj}.) A shift? (\textbf{\mr{पाळी}}, \emph{pāḷī}.) To get done? (\textbf{\mr{उरकणे}}, \emph{urakṇe}.) To struggle? (\textbf{\mr{धडपडणे}}, \emph{dhaḍpaḍṇe}.) An engineer? (\textbf{\mr{अभियंता}}, \emph{abhiyantā}.) A consultant? (\textbf{\mr{सल्लागार}}, \emph{sallāgār}.) A translator? (\textbf{\mr{अनुवादक}}, \emph{anuvādak}.) An accountant? (\textbf{\mr{लेखापाल}}, \emph{lekhāpāl}.) To wrap up? (\textbf{\mr{आटोपणे}}, \emph{āṭopṇe}.) A computer? (\textbf{\mr{संगणक}}, \emph{saṅgaṇak}.) A mobile phone? (\textbf{\mr{भ्रमणध्वनी}}, \emph{bhramaṇdhvanī}.) A telephone? (\textbf{\mr{दूरध्वनी}}, \emph{dūrdhvanī}.) A video clip? (\textbf{\mr{चित्रफीत}}, \emph{citraphīt}.) To use? (\textbf{\mr{वापरणे}}, \emph{vāparṇe}.) A password? (\textbf{\mr{संकेतशब्द}}, \emph{saṅketśabda}.) The internet? (\textbf{\mr{आंतरजाल}}, \emph{āntarjāl}.) A message? (\textbf{\mr{संदेश}}, \emph{sandeś}.) A link? (\textbf{\mr{दुवा}}, \emph{duvā}.) To store? (\textbf{\mr{साठवणे}}, \emph{sāṭhavṇe}.) A machine? (\textbf{\mr{यंत्र}}, \emph{yantra}.) A device? (\textbf{\mr{उपकरण}}, \emph{upakaraṇ}.) A robot? (\textbf{\mr{यंत्रमानव}}, \emph{yantramānav}.) Automatic? (\textbf{\mr{स्वयंचलित}}, \emph{svayañcalit}.) To spray? (\textbf{\mr{फवारणे}}, \emph{phavārṇe}.) A torch? (\textbf{\mr{विजेरी}}, \emph{vijerī}.) A chandelier? (\textbf{\mr{झुंबर}}, \emph{jhumbar}.) Daylight? (\textbf{\mr{उजेड}}, \emph{ujeḍ}.) Darkness? (\textbf{\mr{अंधार}}, \emph{andhār}.) To light up? (\textbf{\mr{उजळणे}}, \emph{ujaḷṇe}.)
\end{tcolorbox}
